\documentclass[./测控三级项目.tex]{subfiles} % 使用 subfiles 文档类，指定主文档

\begin{document}
	\subsection{题目重述}
		\begin{compactenum}
			\item 选用适当参数，完成10倍同相放大电路的Multisim仿真，要求完成题目给定测试。
			\item 焊装电路，做如下测试，并分别叙述与相同输入下仿真波形的差异。
			\item 通过拆换元件，在同相放大电路区域构建电压跟随器。
			\item 恢复10倍同相放大电路。
			\item 说明更换电路的思路和耗损元件，描述相同电路使用不同芯片有无差异，相同芯片构成的不同电路性能参数的异同。
		\end{compactenum}
	\subsection{题目分析与实施方案}
		\subsubsection{第一问}
			第一问需要构建一个同相放大电路，并且实现10倍的放大倍数。然后对这个电路在不同输入信号下的放大效果进行测试，输入信号包括不同频率的正弦波、特定占空比的矩形波和三角波，通过仿真来观察输出信号的特性。
		\subsubsection{第二问}
			第二问依据之前设计的10倍同相放大电路原理图，在PCB上进行电路焊装。并对焊装后的 10 倍同相放大电路进行多方面测试，并对比与仿真波形的差异。对不同峰值和频率的正弦波、矩形波、三角波信号放大测试，能全面了解电路在实际焊装后的性能表现，包括频率响应、压摆率限制以及阶跃响应等特性，进而分析实际电路与理想仿真情况的偏差原因。
		\subsubsection{第三问}
			第三问主要对同相放大电路和电压跟随器电路原理图进行对比，深入理解其与同相放大电路在结构和原理上的内在联系，随后对元件拆换重新构建电压跟随器。并利用构建好的电压跟随器，通过调整信号频率来测量 OP07 的带宽，并与手册参数以及 NE5532 跟随器进行对比，从而全面评估 OP07 在该应用场景下的频率特性以及与其他常用芯片的性能差异。
		\subsubsection{第四问}
			第四问需要按照原有放大电路将10倍同相放大电路进行还原。
		\subsubsection{第五问}
			第五问使用控制变量法，对比不同运放芯片搭建同一电路结构（电压跟随器）时，芯片静态参数和动态参数的异同；对比统一芯片搭建不同电路结构（10倍同相放大电路和跟随器电路），芯片静态参数和动态参数之间的异同。
	
	\subsection{理论分析与仿真}
		\subsubsection{第一问}
			图\ref{fig:同相放大电路}是一个基本的同相放大电路。
			\begin{figure}[h]
				\centering
				\includegraphics[width=0.45\textwidth]{同相放大电路}
				\caption{同相放大电路原理图}
				\label{fig:同相放大电路}
			\end{figure}\par
			由于差分输入电阻$r_{id}$较大，一般认为同相和反相输入端直接没有电流流经，直接忽略该电流。于是，可依据电路结构，写出同相输入端$v_P$和反相输入端$v_N$的电压公式
			\begin{align}
				v_P&=v_I \\
				\frac{v_N-0}{R_1}&=\frac{v_O-v_N}{R_2}
			\end{align}\par
			结合等式$u_O=A_{od}u_I=A_{od}(v_P-v_N)$，联立可得到
			\begin{align}
				v_O=A_{od}(v_I-\frac{1}{1+R_1/R_2}v_O)
			\end{align}\par
			定义实际放大倍数$A$为输出电压$v_O$和输入电压$v_I$之比，则有
			\begin{align}\label{eq:实际同相放大倍数}
				A=\frac{v_O}{v_I}=(1+\frac{R_2}{R_1})\frac{1}{1+(1+R_2/R_1)/A_{od}}
			\end{align}\par
			在式\ref{eq:实际同相放大倍数}中令$A_{od}\rightarrow \infty$就得到理想的放大倍数$A_{ideal}$:
			\begin{align}\label{eq:理想同相放大倍数}
				A_{ideal}=\lim_{A_{od}\rightarrow \infty}A=1+\frac{R_2}{R_1}
			\end{align}\par 
			根据理想放大倍数计算公式式\ref{eq:理想同相放大倍数}，为了得到10倍同相放大电路，只需要满足等式$R_2=9R_1$即可。\par
			以上理论完备，在Multisim软件中进行仿真实验。\par
			根据图\ref{fig:同相放大电路}，搭建仿真电路图，仿真电路使用OP07芯片作为运放芯片，采用$\pm15$V正负电压作为电源供电，反馈电阻网络选择$R_1=1$k$\Omega$，$R_2=9$k$\Omega$电阻将信号发生器接入同相放大电路输入端，示波器通道A接入同相放大电路输入端，示波器通道B接入同相放大电路输出端，其搭建示意图如图\ref{fig:同相放大电路仿真原理图}所示。\par
			\begin{figure}[H]
				\centering
				\includegraphics[width=0.25\textwidth]{./3.第三题示波器现象/同相放大电路仿真原理图}
				\caption{同相放大电路仿真原理图}
				\label{fig:同相放大电路仿真原理图}
			\end{figure}\par
			该仿真电路图测得对应的仿真正弦波形、矩形波形和三角波形如图\ref{fig:仿真同相10倍放大器放大信号}所示。
			\begin{figure}[H]
				\centering
				\begin{subfigure}{0.4\linewidth}
					\centering
					\includegraphics[width=0.75\linewidth]{./3.第三题示波器现象/示波器1}
					\caption{仿真20mV、100Hz正弦信号同相放大10倍}
					\label{fig:仿真同相10倍放大器放大20mV、100Hz正弦信号}
				\end{subfigure}
				\begin{subfigure}{0.4\linewidth}
					\centering
					\includegraphics[width=0.75\linewidth]{./3.第三题示波器现象/示波器2}
					\caption{仿真20mV、1MHz正弦信号同相放大10倍}
					\label{fig:仿真同相10倍放大器放大20mV、1MHz正弦信号}
				\end{subfigure}
				\begin{subfigure}{0.4\linewidth}
					\centering
					\includegraphics[width=0.75\linewidth]{./3.第三题示波器现象/示波器3}
					\caption{仿真2V、100kHz正弦信号同相放大10倍}
					\label{fig:仿真同相10倍放大器放大2V、100kHz正弦信号}
				\end{subfigure}
				\begin{subfigure}{0.4\linewidth}
					\centering
					\includegraphics[width=0.75\linewidth]{./3.第三题示波器现象/示波器4}
					\caption{仿真20mV、100Hz方波信号同相放大10倍}
					\label{fig:仿真同相10倍放大器放大20mV、100Hz、占空比70方波信号}
				\end{subfigure}
				\centering
				\begin{subfigure}{0.4\linewidth}
					\centering
					\includegraphics[width=0.75\linewidth]{./3.第三题示波器现象/示波器5}
					\caption{仿真20mV、100Hz三角波同相放大10倍}
					\label{fig:仿真同相10倍放大器放大20mV、100Hz三角波信号}
				\end{subfigure}
				\caption{仿真同相10倍放大器放大信号}
				\label{fig:仿真同相10倍放大器放大信号}
			\end{figure}\par
			观察图\ref{fig:仿真同相10倍放大器放大20mV、100Hz正弦信号}、\ref{fig:仿真同相10倍放大器放大20mV、100Hz、占空比70方波信号}和\ref{fig:仿真同相10倍放大器放大20mV、100Hz三角波信号}的波形跟随情况，可以发现当峰值都较小（为20mV）且频率较低时，运算放大器具有良好的信号等比放大效果，且放大倍率稳定，没有波形失真问题。\par
			观察图\ref{fig:仿真同相10倍放大器放大20mV、1MHz正弦信号}的波形，可以发现，当输入信号的频率升高时，受到运放芯片压摆率的影响，信号的增益已然变小，从原来的10倍变为了0.5倍，反而使输出信号更加小。同时，输出信号也受到芯片自身动态频率响应问题的影响，造成了巨大的相移，约为$45^\circ$。\par
			观察图\ref{fig:仿真同相10倍放大器放大2V、100kHz正弦信号}的波形跟随情况，可以发现，当输入信号幅值较高时，受到压摆率的影响，运放无法及时的放大信号，造成放大增益的改变。
		\subsubsection{第三问}
			对比图\ref{fig:电压跟随器电路}的电压跟随器原理图和图\ref{fig:同相放大电路}的同相放大电路的原理图，可以发现如果将图\ref{fig:同相放大电路}的同相放大电路的原理图中电阻阻值设为$R_1=+\infty$，$R_2=0\Omega$，则变为了电压跟随器电路。运用公式\ref{eq:理想同相放大倍数}也可以计算得到此时的理想放大倍数$A_{idea}=1+\frac{R_2}{R_1}=1+\frac{0}{+\infty}=1$，理想增益为1，因此电路确认为电压跟随器。\par
	
	\subsection{实际电路测试结果与数据分析}
		\subsubsection{第二问}
			根据等式$R_2=9R_1$，若其中$R_2$选用阻值为$=10$k$\Omega$的电阻，则可以得到$R_1$的阻值为$R_1=1.\dot{1}$k$\Omega$，与标准电阻值表比对后，实际选用电阻阻值为$R=1.1$k$\Omega$(Standeard Value,$1\%$)的电阻。\par
			实际电路板除了以上$R_1$、$R_2$电阻以外还存在输入平衡电阻$R_3$，该电阻是为了维持运算放大器输入级差分电路的对称性，因此该电阻等于负输入端的对地等效电阻，即$R_3=\frac{1}{\frac{1}{R_1}+\frac{1}{R_2}}=\frac{R_1R_2}{R_1+R_2}=1$k$\Omega$。\par
			焊接后，经检查分立元件位置端正、放置平整、焊点表面光滑、无凸起毛刺、无虚焊，周围无助焊剂残留；多引脚器件无短路，焊点均匀一致。符合测试要求。\par
			按第一题的方案连接电源，将信号发生器信号加载在同相放大器信号输入端，将示波器通道CH1接入同相放大器输入端，将示波器通道CH2接入同相放大器输出端。添加除法算子计算输出通道CH2除以输入通道CH1的增益值。\par
			将电路连接完毕后，先打开线性电源输出开关，随后根据《测控仪器电路项目作业指导书》中要求，加载对应正弦信号和方波信号。\par
			同相10倍放大器实验实测对应的仿真正弦波形、矩形波形和三角波形如图\ref{fig:实际同相10倍放大器放大信号}所示。
			\begin{figure}[H]
				\centering
				\begin{subfigure}{0.4\linewidth}
					\centering
					\includegraphics[width=0.75\linewidth]{./3.第三题示波器现象/UpoPrtsc-001}
					\caption{实际20mV、100Hz正弦信号同相放大10倍}
					\label{fig:实际同相10倍放大器放大20mV、100Hz正弦信号}
				\end{subfigure}
				\begin{subfigure}{0.4\linewidth}
					\centering
					\includegraphics[width=0.75\linewidth]{./3.第三题示波器现象/UpoPrtsc-002}
					\caption{实际20mV、1MHz正弦信号同相放大10倍}
					\label{fig:实际同相10倍放大器放大20mV、1MHz正弦信号}
				\end{subfigure}
				\begin{subfigure}{0.4\linewidth}
					\centering
					\includegraphics[width=0.75\linewidth]{./3.第三题示波器现象/UpoPrtsc-003}
					\caption{实际2V、100kHz正弦信号同相放大10倍}
					\label{fig:实际同相10倍放大器放大2V、100kHz正弦信号}
				\end{subfigure}
				\begin{subfigure}{0.4\linewidth}
					\centering
					\includegraphics[width=0.75\linewidth]{./3.第三题示波器现象/UpoPrtsc-004}
					\caption{实际20mV、100Hz方波信号同相放大10倍}
					\label{fig:实际同相10倍放大器放大20mV、100Hz、占空比70方波信号}
				\end{subfigure}
			\end{figure}
			\begin{figure}[H]\ContinuedFloat
				\centering
				\begin{subfigure}{0.4\linewidth}
					\centering
					\includegraphics[width=0.75\linewidth]{./3.第三题示波器现象/UpoPrtsc-005}
					\caption{实际20mV、100Hz三角波同相放大10倍}
					\label{fig:实际同相10倍放大器放大20mV、100Hz三角波信号}
				\end{subfigure}
				\caption{实际同相10倍放大器放大信号}
				\label{fig:实际同相10倍放大器放大信号}
			\end{figure}\par
			图\ref{fig:实际同相10倍放大器放大20mV、100Hz正弦信号}的信号相比于图\ref{fig:仿真同相10倍放大器放大20mV、100Hz正弦信号}图像中完全重叠的信号，在零点还有较多毛刺，这是噪声信号和运算放大器动态频率响应共同作用导致的问题。这说明实际运放电路搭建中，会随着信号的处理，不可避免的引入噪声和相移。电路设计者需注意对这些随机量的处理。\par
			图\ref{fig:实际同相10倍放大器放大20mV、1MHz正弦信号}和图\ref{fig:仿真同相10倍放大器放大20mV、1MHz正弦信号}在输出信号的大致现象上是一致的，输出信号频率超过运放在该增益上的截止频率后会随着信号频率的进一步增高而增益减小，并产生相移现象。这说明了运放的增益带宽积会对高频信号的放大增益产生限制。但是，另一面，仿真图像的增益和相移，与实际测量图像有一定的差距，这意味着Multisim仿真平台对于运算放大器的建模仿真是不完全和实际重合的，仿真只能作为电路设计的一个参考。\par
			图\ref{fig:实际同相10倍放大器放大2V、100kHz正弦信号}和图\ref{fig:仿真同相10倍放大器放大2V、100kHz正弦信号}的图像较为接近。对于频率较高且振幅较大的信号，运放会受到其压摆率的影响，对于输入信号有较大的相移，并伴随着压降。\par
			图\ref{fig:实际同相10倍放大器放大20mV、100Hz、占空比70方波信号}和图\ref{fig:仿真同相10倍放大器放大20mV、100Hz、占空比70方波信号}的图像基本一致。将图像放大来看，信号在上升沿、下降沿处是个阶跃响应，运放芯片会随着信号幅值的改变产生一定的量的超调。幅值越大，超调量越大。同时也调试发现，随着频率信号的升高，阶跃响应的上升时间会加长，这是因为对于同一芯片的增益带宽积是一定的，随着信号频率的增大，运放的压摆率是逐渐减小的，使得上升时间的延长。\par
			图\ref{fig:实际同相10倍放大器放大20mV、100Hz三角波信号}和图\ref{fig:仿真同相10倍放大器放大20mV、100Hz三角波信号}的图像基本一致。
			
		\subsubsection{第三问}
			根据理论分析与仿真中对同相放大电路和电压跟随器电路的比较，得出电路修改方案。根据方案，拆除原本焊接的电阻，并短接$R_2$和$R_3$的焊盘。\par
			修改电路后，按第一题的方案连接电源，将信号发生器信号加载在电压跟随器信号输入端，将示波器通道CH1接入电压跟随器输入端，将示波器通道CH2接入电压跟随器输出端。\par
			将电路连接完毕后，先打开线性电源输出开关，随后不断调整信号发生器输出的正弦波频率，直到示波器通道CH2中测得的输出信号幅度降至输出信号的$\frac{\sqrt{2}}{2}$时，即增益为$\frac{\sqrt{2}}{2}$时，停止测量，此时正弦信号的频率即为该运放的带宽。\par
			\begin{figure}[H]
				\centering
				\begin{subfigure}{0.4\linewidth}
					\centering
					\includegraphics[width=0.75\linewidth]{./3.第三题示波器现象/UpoPrtsc-007}
					\caption{OP07CP截止频率处输出信号截屏}
					\label{fig:OP07CP截止频率处输出信号截屏}
				\end{subfigure}
				\begin{subfigure}{0.4\linewidth}
					\centering
					\includegraphics[width=0.75\linewidth]{./3.第三题示波器现象/UpoPrtsc-008}
					\caption{NE5532截止频率处输出信号截屏}
					\label{fig:NE5532截止频率处输出信号截屏}
				\end{subfigure}
				\caption{不同运放截止频率处输出信号截屏}
				\label{fig:不同运放截止频率处输出信号截屏}
			\end{figure}\par
			经过测量，当OP07CP和NE5532分别在频率为850kHz和6MHz时，信号幅值约为0.035mV，是输入信号0.05mV的$\frac{\sqrt{2}}{2}$倍，其输入图像如图\ref{fig:不同运放截止频率处输出信号截屏}所示。\par
			查阅数据手册可得，OP07CP的开环增益带宽典型值为0.6MHz，NE5532的开环增益带宽典型值为10MHz。
		\subsubsection{第四问}
			将卸下的元件重新恢复到PCB上。焊接后，经检查分立元件位置端正、放置平整、焊点表面光滑、无凸起毛刺、无虚焊，周围无助焊剂残留；多引脚器件无短路，焊点均匀一致。
		
		\subsubsection{第五问}
			跟换电路的思路是根据同相放大电路的放大倍数公式，和电压跟随器增益为1的特点，产生使反馈电阻为0，并让输入电阻趋于无穷大的手段得到两个电路的跟换效果。\par
			为防止浪费元器件，我们将元件卸下后收齐起来，等测试完芯片带宽后又继续按照电阻上的丝印确认电阻阻值还原回10倍同相放大电路。\par
			对于相同电路，不同芯片在低频区间的区别不大，主要实在高频信号的处理上会有不同的信号频率响应。由于OP07CP的带宽远小于NE5532，因此在高频信号如1MHz的正弦信号，可以明显的看见相比于NE5532，OP07CP的跟随信号有增益降低，相位偏移的问题。
	\subsection{结论与收获}
		\subsubsection{结论}
			\paragraph{\qquad 10 倍同相放大电路 Multisim 仿真结论}
			对峰值 20mV、100Hz 正弦信号放大，输出波形为正弦，幅值约为 200mV，相位与输入相同，实现了 10 倍放大，波形无明显失真。\par
			对峰值 20mV、1MHz 正弦信号放大，输出幅值小于理论值，波形有一定失真，表明运放带宽对高频信号放大有影响，高频信号放大倍数下降。\par
			对峰值 2V、100KHz 正弦信号放大，输出波形正弦，能达到一定放大倍数，但与理论 10 倍放大有偏差，可能受元件参数非理想等因素影响。\par
			对峰值 20mV、占空比 70\%、频率 100Hz矩形波信号放大，输出矩形波，上升沿和下降沿较陡，波形形状保持较好，实现了 10 倍放大。\par
			对峰值 20mV、100Hz 三角波信号放大，输出三角波，斜率增大，实现 10 倍放大，波形无明显失真。
			\paragraph{\qquad 焊装电路测试结论}
			内部失调调零可减小输出失调电压，使输出更稳定准确。\par
			对峰值 20mV、100Hz 正弦信号放大，实测波形与仿真基本一致，验证了电路实际性能与仿真的相符性。\par
			对峰值 20mV、1MHz 正弦信号放大，与仿真波形类似，高频时输出波形有失真趋势，进一步证明频率对输出波形影响。\par
			对峰值 2V、100KHz 正弦信号放大，实测与仿真有误差，可能因实际元件参数、电路寄生参数等影响，且摆率影响使输出波形顶部和底部有削波现象。\par
			对峰值 20mV、占空比 70\%、频率 100Hz 矩形波信号放大，实测波形与仿真接近，阶跃响应良好，输出能快速跟随输入变化。\par
			对峰值 20mV、100Hz 三角波信号放大，实测与仿真基本一致，电路对三角波信号放大性能稳定。
			\paragraph{\qquad 电压跟随器相关结论}
			电压跟随器结构为反馈电阻为零，输出接反相输入端，信号源接同相输入端，其原理基于运放特性使输出等于输入，具有高输入阻抗和低输出阻抗，可作缓冲器。\par
			通过短路反馈电阻和断开输入电阻将同相放大电路转换为电压跟随器，反之可恢复。转换过程中可能损坏电阻、电容及电路板线路和焊盘。\par
			测量 OP07 带宽约为 600kHz，与手册参数相符且远小于 NE5532 跟随器带宽（10MHz），OP07 适用于低频高精度应用，NE5532 适用于高速信号处理领域。
			\paragraph{\qquad 相同电路不同芯片及相同芯片不同电路结论}
			不同芯片（如 OP07 和 NE5532）在带宽、压摆率、输入输出特性、噪声特性等性能上有差异，适用于不同应用场景。\par
			相同芯片构成的同相放大电路和电压跟随器，固有特性（如失调电压、电源要求）相同，但放大倍数、带宽、输入输出阻抗等性能参数不同。
		\subsubsection{小组成员收获}
			\paragraph{\qquad 理论分析层面}
			深入理解同相放大电路原理，熟练掌握放大倍数与电阻关系推导，能根据需求设计计算参数，明确不同输入信号下电路响应理论依据。透彻理解电压跟随器结构原理及其与同相放大电路关系，知晓其为同相放大电路特殊形式及在信号处理中的不同作用。掌握电路性能参数（如带宽、压摆率、输入输出阻抗等）对信号放大的影响机制，学会分析信号频率、幅值与电路性能的相互关系，能从理论上解释输出波形变化原因。熟悉不同芯片（如 OP07 和 NE5532）的性能特点及适用场景，理解芯片性能差异对电路整体性能的影响，为电路设计中芯片选型提供理论支持。
			\paragraph{\qquad 动手实践方面}
			熟练运用 Multisim 进行电路仿真，学会设置不同输入信号参数，观察分析输出波形，通过仿真预先了解电路性能。掌握电路焊装技能，包括元件选型、布局、焊接及电路调试方法，能准确构建实际电路并使其正常工作。学会使用示波器、信号发生器等仪器进行电路测试，能准确测量波形参数，根据测试结果判断电路性能优劣。通过拆换元件构建不同电路（如电压跟随器），提高电路改造和优化能力，积累实际操作经验，学会解决电路搭建和调试中遇到的问题（如元件损坏、参数偏差等）。
			\paragraph{\qquad 精神层面}
			培养团队协作精神，小组成员分工合作完成各项任务，如电路设计、仿真、焊装、测试等，通过沟通协作解决问题，提高团队工作效率。增强耐心和细心，在电路设计计算、元件选型、焊装调试等过程中，任何小失误都可能影响结果，需耐心细致操作和检查，逐渐养成严谨工作态度。培养勇于探索和创新精神，在遇到电路性能不佳等问题时，积极尝试不同解决方案，如调整元件参数、更换芯片等，不断探索优化电路性能的方法，提升创新思维能力。面对多次测试失败和复杂问题时，不气馁不放弃，坚持寻找解决办法，培养坚韧不拔的毅力和克服困难的决心。
			
	
	
\end{document}